\documentclass[11pt,a4paper]{article}
\usepackage[margin=25mm]{geometry}
\usepackage[T1]{fontenc}
\usepackage{lmodern}
\usepackage{microtype}
\usepackage{xcolor}
\usepackage[hidelinks]{hyperref}
\usepackage{parskip}
\pagestyle{empty}

\definecolor{ink}{HTML}{25324A}
\definecolor{muted}{HTML}{5C667A}

\begin{document}

{\color{ink}\LARGE\bfseries Demographics accuracy: short source}

\vspace{2mm}
{\color{muted}\small Guided reading material - 11 August 2026}

\vspace{7mm}
\textbf{Read the adapted extract below.} It preserves the study's method, central projection, main regional explanation, and uncertainty. It is not a quotation from the paper.

\vspace{4mm}
Gerland and colleagues analysed United Nations population data available up to 2012. Instead of presenting only a few fixed scenarios, they used a Bayesian probabilistic method to estimate a range of possible population outcomes.

Their model indicated an 80 percent probability that the world population, then about 7.2 billion, would rise to between 9.6 and 12.3 billion by 2100. The authors therefore argued that global population growth was unlikely to stop during the twenty-first century.

Much of the projected increase was expected to occur in Africa. The researchers linked this pattern partly to relatively high fertility and to a recent slowdown in the rate at which fertility was declining.

The study also projected a substantial fall in the ratio of working-age people to older people across countries. These statements are probabilistic projections based on data and model assumptions: they describe likely global and regional futures, not a certain result for every individual country.

\vfill
{\color{muted}\small
Adapted from: Gerland, P. et al. (2014), ``World population stabilization unlikely this century,'' \textit{Science}, 346(6206), 234-237. DOI: \href{https://doi.org/10.1126/science.1257469}{10.1126/science.1257469}. The adaptation is for guided study and should not be cited as a verbatim extract.
}

\end{document}
